\chapter*{Abstract}
\addcontentsline{toc}{chapter}{Abstract}

Data-driven forecasting models select their predictors by statistical association, yet in the tropical atmosphere most associations are inflated by autocorrelation, shared seasonal cycles and common drivers. This thesis asks whether a \emph{validated} graph of lagged climate dependencies, found by causal discovery and frozen before any test year is seen, can tell a forecasting model which source regions to use. The setting is daily outgoing longwave radiation (OLR), a proxy for deep tropical convection, forecast for 50 tropical regions centred on the Indo-Pacific during the boreal summer intraseasonal oscillation (BSISO) season.

Phase~1, the primary contribution, represents the climate system as 300 time-series nodes (50 regions $\times$ six variables: SST, 850\,hPa geopotential height, 200 and 850\,hPa zonal wind, precipitable water and OLR) over 1979--2022, with 1979--2012 reserved for discovery and training and 2013--2022 held out. A BSISO index reproduced from the project's own data (leading modes explaining 11.2\% and 7.7\% of variance, close to published values) defines an activity-based season that restricts target days only. PCMCI with partial-correlation tests and Benjamini--Hochberg false-discovery-rate control is run over 18{,}000 lagged candidates per OLR target (lags 1--60 days) using a vectorised implementation verified link-for-link against Tigramite. The resulting graph of 536 links is passed through four independent checks: phase-randomised decoys (0 of 250 decoys selected), a physical propagation-speed mask (3 of 536 links rejected), data-driven top-$k$ selection, and persistence at doubled spatial resolution (mean persistence 0.62; 33 of 48 scored targets at or above 0.5). The validated graph is sparse, short-lag and largely regional: a median outside-link lag of one day and about two source regions per target.

The frozen graph is then used as a source-selection prior for an invariance-trained causal splitter and per-region expert within a Mixture-of-Experts forecasting framework. In the full 50-region candidate pool the graph-restricted splitter becomes discriminative (informative splitters rise from 0 to 24 of 50 regions at a 7-day lead and from 0 to 36 at a 1-day lead). Seven-day forecast skill is unchanged relative to the no-prior model (mean skill against persistence $+0.416$ versus $+0.415$), whereas at a 1-day lead skill improves in 32 of 50 regions (mean skill $-1.33$ versus $-1.53$; mean $R^2$ against climatology $0.745$ versus $0.720$). Because the graph links act on time scales of one to two days, the evidence locates the graph's value at short leads, and attributing the gain to the graph's choice of regions, rather than to pool size, requires the matched random- and correlation-prior controls defined in the experimental design. A sliding-window analysis shows that changes in the causal structure do not align with ENSO phase changes at one- and three-season window lengths.

Phase~2 extends the framework to hourly, kilometre-scale near-surface forecasting over India, combining a shared backbone, causally gated regional experts and drift-triggered expert adaptation, and specifies the data, baselines and shift-aware evaluation it requires. The thesis contributes a reproducible build--validate--freeze--use protocol for causal graphs in forecasting, a validated graph with full diagnostics, and an evaluation design that separates a graph's credibility from its predictive usefulness.

\vspace{0.4cm}
\noindent\textbf{Keywords:} causal discovery, PCMCI, BSISO, outgoing longwave radiation, surrogate testing, mixture of experts, graph prior, subseasonal forecasting, concept drift.
